\section{Classes}

% [CORRECTED] was: ``so $V$ has no $\in$-minimal member'' --- $\emptyset \in V$ is $\in$-minimal in $V$, so Foundation holds for $V$ itself; the set it fails for is $\set{V}$, whose only member $V$ satisfies $V \in V$. The same repair is made where the argument is restated below.
Recall Russell's Paradox: there is no set of all sets. One reason for this (which relies on the Axiom of Foundation) is that if $V$ is the set of all sets, then $V \in V$, so $\set{V}$ has no $\in$-minimal member, which contradicts the Axiom of Foundation. Without the Axiom of Foundation, if $V$ is the set of all sets, then if we set
\begin{align*}
    W = \setst{x \in V}{x \notin x}
\end{align*}
(which is a set by the Axiom of Comprehension), then if $W \in W$, then $W \notin W$, and if $W \notin W$, then $W \in W$, which is a contradiction, meaning there can be no set of all sets.

Nevertheless, we really do want to write things like
\begin{align*}
    V = \setst{x}{x \text{ is a set}} = \setst{x}{x = x}
\end{align*}
We would also like to write things like
\begin{align*}
% [CORRECTED] was: two class terms on consecutive source lines of one align* row, typesetting as $\OR = \setst{\alpha}{\ldots} \setst{V_{\alpha}}{\ldots}$ with nothing between them; they are two separate examples of things we would like to write
    \OR = \setst{\alpha}{\alpha \text{ is an ordinal}} \qquad \text{and} \qquad \setst{V_{\alpha}}{\alpha \in \OR}
\end{align*}
The only problem is, none of these are sets. Indeed, the only guarantee of the Axiom of Comprehension is that
\begin{align*}
    \setst{x \in A}{\varphi\of{x, \bar{y}}}
\end{align*}
is a set; there is no guarantee whatsoever that
\begin{align*}
    \setst{x}{\varphi\of{x, \bar{y}}}
\end{align*}
is a set. It sometimes is, but it isn't always.

% [CORRECTED] was: ``meaning it would have no $\in$-minimal element'' --- same fault as at the top of the section: $\emptyset$ is $\in$-minimal in $V$; it is $\set{V}$ that Foundation fails for
There are two ways you can prove that $V = \setst{x}{x = x}$ is not a set. The first is that if it were, then $V \in V$ would hold, meaning $\set{V}$ would have no $\in$-minimal element, contradicting foundation. The second is that if it were, $W = \setst{x \in V}{x \notin x}$ would be a valid set, by comprehension. Then, $W \in W \lr W \notin W$, a contradiction.

There are two perspectives we can take. The first is the \textbf{external perspective}: this is the perspective from which you would look at models of ZFC. There is also the internal perspective, where you are already inside a model of ZFC. Now this internal perspective also has a notion of logic and set theory inside of it. So there's formalisations of set theory inside models of set theory, and these formalisations have their own notions of models of set theory, and so it goes.

At some point, we will study these things in detail.

For the time being, we will use the term \textbf{class} to describe something we might want to express as a set, even if it cannot actually be a set for cardinality reasons.

Note that we would want sets to be classes. Indeed, anything of the form $\setst{x}{\varphi(x)}$, with $\varphi(x)$ a formula in the language of set theory, should be a class. Unfortunately, this is not generous enough: we want, for each set $A$, that
\begin{align*}
    A = \setst{x}{x \in A} = \setst{x}{\varphi\of{x, A}}
\end{align*}
to be included as well. In other words, we want to have formulae $\varphi$ that not only allow us to substitute free variables with input sets $x$ but also allow us to substitute more free variables with parameter sets $A$.

Formulae and parameters determine the class, but the converse is not true. In other words, if
\begin{align*}
    \forall x \parenth{\varphi(x) \lr \varphi'(x)}
\end{align*}
then
\begin{align*}
    \setst{x}{\varphi(x)} = \setst{x}{\varphi'(x)}
\end{align*}
For example, we know that
\begin{align*}
    \varphi(x) \lr \parenth{\varphi(x) \land x = x}
\end{align*}
So you can define the same class using more than one formula.

\subsection{The Class of Ordinals}

\begin{boxdefinition}[The Class of Ordinals]
    Define
    \begin{align*}
        \OR := \setst{\alpha}{\alpha \text{ is an ordinal}}
    \end{align*}
    We can show that $\OR$ is a class.
\end{boxdefinition}

We can then say that the `mapping'
\begin{align*}
    \alpha \mapsto V_{\alpha}
\end{align*}
is a \textbf{class function}: it takes elements of a class and maps them to elements of another class.

We can represent this mapping as
\begin{align*}
    \setst{\parenth{\alpha, V_{\alpha}}}{\alpha \in \OR}
\end{align*}
Now we can ask ourselves the following question: for which formulae $\varphi$ is this set equal to $\setst{x}{\varphi(x)}$?

Let's give two such formulae based on the following.

\begin{boxtheorem}
    $\forall \gamma \in \OR$, $\forall S \in V$, TFAE:
    \begin{enumerate}
        \item $S = V_{\gamma}$
        \item $\exists g$ a function with the properties that
        \begin{enumerate}[label = (\alph*)]
            \item The domain of $g$ is $\gamma + 1$
            \item $g(0) = \emptyset$
            \item $\forall \alpha < \gamma$, $g\of{\alpha + 1} = \Powset\of{g\of{\alpha}}$
            \item $\forall \beta \leq \gamma$ with $\beta$ a limit ordinal, $g\of{\beta} = \bigcup_{\alpha < \beta} g(\alpha)$
            \item $g(\gamma) = S$
        \end{enumerate}
        \item Every function $g$ exhibiting properties (a)-(d) above also exhibits property (e).
    \end{enumerate}
\end{boxtheorem}

\subsection{Well-Foundedness}

We now define the class of well-founded sets.

\begin{boxdefinition}[The Class of Well-Founded Sets]
    Define
    \begin{align*}
        \WF := \bigcup_{\alpha \in \OR} V_{\alpha}
    \end{align*}
    That is, $\WF$ consists of all $x$ such that there is some ordinal $\alpha$ with $x \in V_{\alpha}$.
\end{boxdefinition}

$\WF$ is a class.

We will eventually see that the axiom of foundation is equivalent to saying that $V = \WF$ (under $\ZF$ without foundation).

\begin{boxproposition}
    For all $\gamma \in \OR$,
    \begin{enumerate}[label = $\parenth{\arabic*}_{\gamma}$]
        \item $\forall \alpha \leq \gamma$, $V_{\alpha} \ssq V_{\gamma}$
        \item $V_{\gamma}$ is a transitive set
    \end{enumerate}
\end{boxproposition}
\begin{proof}
    We argue by induction on ordinals $\gamma$.

    For the base step ($\gamma = 0$), both $(1)_0$ and $(2)_0$ are clear.

    For the successor step, if $\gamma = \beta + 1$,
    \begin{enumerate}[label = $\parenth{\arabic*}_{\gamma}$]
        \item Let $\alpha \leq \beta$ and $x \in V_{\alpha}$. By $(1)_{\beta}$, $x \in V_{\beta}$. By $(2)_{\beta}$, $x \ssq V_{\beta}$. So, $x \in \Powset\of{V_{\beta}} = V_{\beta + 1} = V_{\gamma}$.

        \item Let $y \in V_{\gamma} = \Powset\of{V_{\beta}}$. Then, $y \subseteq V_{\beta}$. By $(1)_{\gamma}$, $y \ssq V_{\gamma}$.
    \end{enumerate}

    Finally, for the limit step, ie, if $\gamma$ is a limit ordinal, then $(1)_{\gamma}$ and $(2)_{\gamma}$ both follow from the fact that $\forall \beta < \gamma$, $(1)_{\beta}$ and $(2)_{\beta}$ both hold. (This is because in the limit case, constructing $V_{\gamma}$ is just done by taking unions.)
\end{proof}

\begin{boxdefinition}[Rank]\label{Ch1:Def:Rank}
    Define, for $x \in \WF$, the \textbf{rank} of $x$, denoted $\rank{x}$, to be the least $\alpha \in \OR$ such that $x \in V_{\alpha + 1}$.
\end{boxdefinition}

Indeed, this is equivalent to $\rank{x}$ being the least $\alpha \in \OR$ such that $x \subseteq V_{\alpha}$.

\begin{boxproposition}
    The mapping
    \begin{align*}
        x \mapsto \rank{x}
    \end{align*}
    is a class surjection from $\WF$ to $\OR$.
\end{boxproposition}

We won't prove this, but we'll say it's easy enough to prove and leave it at that.

\begin{boxproposition}
    $\forall \beta \in \OR$,
    \begin{align*}
        V_{\beta} = \setst{x \in \WF}{\rank{x} < \beta}
    \end{align*}
\end{boxproposition}
\begin{proof}[Proof sketch]
    You can just do induction on $\beta$ and apply the definition of the rank (\Cref{Ch1:Def:Rank}).
\end{proof}

It turns out that we can tell whether a set lies in $\WF$ by looking only at its members.

\begin{boxproposition}
    $\forall y \parenth{y \in \WF \lr y \subset \WF}$
\end{boxproposition}
\begin{proof}\hfill
    \begin{description}
        \item[$\to$] Let $\beta = \rank{y}$. Then $y \ssq V_{\beta} \subset \WF$.
        \item[$\leftarrow$] Assume $y \subset \WF$. Define $f : y \to \OR$ by $f(x) = \rank{x} + 1$. We're using replacement to know that $f \in V$.

        Let $\beta = \sup_{x \in y}{f(x)}$. Then, $y \ssq V_{\beta}$. So $y \in \Powset\of{V_{\beta}} = V_{\beta + 1} \subset \WF$.
    \end{description}
\end{proof}

Members have strictly smaller rank than the sets that contain them.

\begin{boxproposition}\label{Ch1:Prop:Rank_Of_Members}
    \begin{align*}
        \forall y \in \WF,\, \forall x \in y,\, \brac{x \in \WF \land \rank{x} < \rank{y}}
    \end{align*}
\end{boxproposition}
\begin{proof}
    We know $x \in y \ssq V_{\rank{y}} \subset \WF$, so $x \in \WF$. This also shows that $\rank{x} + 1 \leq \rank{y}$, so $\rank{x} < \rank{y}$.
\end{proof}

From this, we can compute the rank of a set from the ranks of its members.

\begin{boxproposition}[The Rank Formula]
    \begin{align*}
        \forall y \in \WF, \, \rank{y} = \sup_{x \in y}\of{\rank{x} + 1}
    \end{align*}
\end{boxproposition}
\begin{proof}
    Put $\rho = \rank{y}$ and $\sigma = \sup_{x \in y}\of{\rank{x} + 1}$. By \Cref{Ch1:Prop:Rank_Of_Members}, if $x \in y$, then $\rank{x} < \rank{y} = \rho$, so $\rank{x} + 1 \leq \rho$. So $\sigma \leq \rho$.

    It suffices to show that $y \ssq V_{\sigma}$ to conclude that $\rho \leq \sigma$. Now $V_{\sigma} = \setst{x}{\rank{x} < \sigma}$. Every $x \in y$ has $\rank{x} + 1 \leq \sigma$, by the definition of $\sigma$ as a supremum, and so $\rank{x} < \sigma$. So $y$ is indeed a subset.
\end{proof}

It turns out that the ordinals are their own ranks.

\begin{boxproposition}
    \begin{align*}
        \forall \beta \in \OR, \,
        \brac{\beta \in \WF \land \rank{\beta} = \beta}
    \end{align*}
\end{boxproposition}
\begin{proof}
    We argue by induction on $\beta$. Base case is trivial. In the inductive case, $\beta \ssq \WF$, so by previous prop, $\beta \in \WF$. So,
    \begin{align*}
        \rank{\beta}
        &= \sup_{\alpha < \beta}\of{\rank{\alpha} + 1} \\
        &= \sup_{\alpha < \beta}\of{\alpha + 1} \\
        &= \beta
    \end{align*}
\end{proof}

It follows that each $V_{\beta}$ contains exactly the ordinals below $\beta$.

\begin{boxproposition}
    \begin{align*}
        \forall \beta \in \OR, \,
        \brac{V_{\beta} \cap \OR = \beta}
    \end{align*}
\end{boxproposition}
\begin{proof}
    We know
    \begin{align*}
        V_{\beta} \cap \OR
        &= \setst{x \in \WF}{\rank{x} < \beta} \cap \OR \\
        &= \setst{\alpha \in \OR}{\rank{\alpha} < \beta} \\
        &= \setst{\alpha \in \OR}{\alpha < \beta} \\
        &= \beta
    \end{align*}
\end{proof}

\subsection{Well-Foundedness of Membership}

We now say what it means for a relation on a class to be well-founded. Say $A$ is a class and $R$ is a class binary relation\footnote{ie, a binary relation on elements of the class $A$}. When we write $(A, R)$, we really mean the class structure
\begin{align*}
    \parenth{A, R \cap \parenth{A \times A}}
\end{align*}

\begin{boxdefinition}[Well-Foundedness]
    We say that $(A, R)$ is \textbf{well-founded} if for every nonempty $S \ssq A$, $\exists x \in S$ such that $\forall y \in S$, $\neg \parenth{y \, R \, x}$.

    We say that $(A, R)$ is \textbf{ill-founded} if it is not well-founded.
\end{boxdefinition}

Membership on a well-founded set is, reassuringly, well-founded.

\begin{boxproposition}\label{Ch1:Prop:WF_Well_Founded}
    $\forall A \in \WF$, $\parenth{A, \in}$ is well-founded.
\end{boxproposition}
\begin{proof}
    Let $S \ssq A$, $S \ne \emptyset$. Denote
    \begin{align*}
        \alpha = \min_{x \in S} \rank{x}
    \end{align*}
    % [FILLED] the lecture broke off after naming $\alpha$; the rest of the argument is supplied here
    This makes sense: every $x \in S$ is a member of $A \in \WF$, so $x \in \WF$ by \Cref{Ch1:Prop:Rank_Of_Members} and has a rank, and $\OR$ is well-ordered. Pick $x \in S$ with $\rank{x} = \alpha$. If some $y \in S$ had $y \in x$, then $\rank{y} < \rank{x} = \alpha$ by \Cref{Ch1:Prop:Rank_Of_Members} again, contradicting the minimality of $\alpha$. So $x$ is $\in$-minimal in $S$.
\end{proof}

Before going further, we give a name to the smallest transitive set containing a given set; that there is one is the content of a lemma below.

\begin{boxdefinition}[Transitive Closure]
    Given a set $Y$, the \textbf{transitive closure} of $Y$, denoted $\trcl{Y}$, is the smallest transitive set $T$ such that $Y \ssq T$. In particular, $\trcl{\set{x}}$ is the smallest transitive set containing $x$ as a member.
\end{boxdefinition}

The converse of \Cref{Ch1:Prop:WF_Well_Founded} fails, and the transitive closure is exactly what the failure is hiding.

\begin{boxnexample}[An Ill-Founded Relation]
    One way Foundation can fail is if we have two sets $x \ne y$ such that $x = \set{y}$ and $y = \set{x}$.
    
    In this case, $(x, \in)$ and $(y, \in)$ are both well-founded (just because $x$ and $y$ are both singletons). But $x, y \notin \WF$, because if they were, then they would \textit{both} be, so they would both have ranks, and since $x \in y$, $\rank{x} < \rank{y}$ and similarly, since $y \in x$, $\rank{y} < \rank{x}$, which is impossible.

    In this example, the transitive closure of $x$ is $\set{x, y}$, which is also the transitive closure of $y$. Moreover, $\parenth{\set{x, y}, \in}$ is ill-founded.
\end{boxnexample}

For transitive sets, on the other hand, there is nothing to hide.

\begin{boxproposition}\label{Ch1:Prop:Transitive_WF}
    For every transitive set $A$, if $(A, \in)$ is well-founded, then $A \in \WF$.
\end{boxproposition}
\begin{proof}
    Fix $A$ a transitive set and assume $(A, \in)$ is well-founded. To see that $A \in \WF$, it suffices to show that $A \subset \WF$.

    Suppose not. Let $S = A \setminus \WF$. Then, $S$ is a nonempty subset of $A$. As $(A, \in)$ is well-founded, there is some $x \in S$ that is $\in$-minimal in $S$. Then $x \notin \WF$, so $x \not\subset \WF$.

    Pick $y \in x \setminus \WF$. Then, $y \in x \in A$. Since $A$ is transitive, $y \in A$. Thus, $y \in A \setminus \WF = S$, and $y \in x$. This contradicts our choice of $x$.
\end{proof}

The transitive closure can be built from below, in $\omega$ steps.

\begin{boxlemma}
    Given a set $Y$, put
    \begin{align*}
        Z_0 &= Y \\
        Z_{n + 1} &= \bigcup Z_n \\
        Z_{\omega} &= \bigcup_{n < \omega} Z_n
    \end{align*}
    Then, $\trcl{Y} = Z_{\omega}$.
\end{boxlemma}

We can now deliver on the promise made earlier in this section: over $\ZFwoF$, Foundation says precisely that $V = \WF$.

\begin{boxproposition}[Assuming $\ZFwoF$]
    TFAE:
    \begin{enumerate}[label = (\arabic*)]
        \item Foundation: every nonempty $S$ has an $\in$-minimal member
        \item For all $A$, $(A, \in)$ is well-founded
        \item $V = \WF$
    \end{enumerate}
\end{boxproposition}
\begin{proof}
    $(1) \iff (2)$ is clear: taking $A = S$ in $(2)$ gives $(1)$, and conversely a nonempty $S \ssq A$ is in particular a nonempty set, so $(1)$ hands us an $\in$-minimal member of $S$.

    $(3) \implies (2)$ follows from \Cref{Ch1:Prop:WF_Well_Founded}.

    Let's now assume $(2)$ and prove $(3)$. $\WF \ssq V$, so we only need $V \ssq \WF$. Let $x \in V$ and let $A = \trcl{\set{x}}$.
    By \Cref{Ch1:Prop:Transitive_WF}, since $A$ is transitive and $(A, \in)$ is well-founded, by our assumption $(2)$, we get that $A \in \WF$. So $x \in A \subset \WF$. Thus, $x \in \WF$.
\end{proof}
